\documentclass[11pt]{article}

\usepackage[margin=1in]{geometry}
\usepackage{amsmath,amssymb}
\usepackage{booktabs}
\usepackage{graphicx}
\usepackage{hyperref}
\usepackage{xcolor}
\usepackage{enumitem}
\usepackage{caption}

\hypersetup{
    colorlinks=true,
    linkcolor=blue!50!black,
    citecolor=blue!50!black,
    urlcolor=blue!50!black
}

\title{\textbf{Adapting Qwen-Image-Layered for Portrait Layer Decomposition}\\
\large An Engineering and Modeling Study with CelebAMask-HQ}
\author{Project Report}
\date{\today}

\begin{document}
\maketitle

\begin{abstract}
This project investigates whether \texttt{Qwen/Qwen-Image-Layered}, a model
designed to decompose images into editable RGBA layers, can be adapted to a
portrait-specific layer decomposition task using CelebAMask-HQ supervision. The
target application is simple to state: given a portrait, produce a small number
of semantically meaningful transparent layers that can be edited independently.
In practice, this became a study of the gap between semantic segmentation targets
and the native decomposition behavior of a pretrained layered image model. We
built an end-to-end Colab pipeline for preprocessing, metadata recovery, sharded
packaging, LoRA training, checkpoint comparison, and numeric evaluation. Stronger
fine-tuning settings on the original five-layer target caused noisy or collapsed
outputs. The later three-layer BG/HAIR/FACE formulation was more stable, even at
200 steps, but still produced only small improvements over the base model. The
final 32-sample evaluation shows a small macro IoU improvement from 0.2461 to
0.2585 at checkpoint 200, driven mostly by the face layer, while layer-health
diagnostics still show high overlap and frequent full-canvas layers. The main
conclusion is that the engineering pipeline is reliable, but target representation
and alpha separation remain the central modeling bottlenecks.
\end{abstract}

\section{Motivation}

Layered image editing is attractive because it turns a flat image into editable
objects. If hair, face, background, text, or other components are isolated in
separate RGBA layers, then downstream edits such as recoloring, moving, deleting,
or replacing content can be applied locally. \texttt{Qwen-Image-Layered} is built
around this idea: it takes an input image and produces multiple RGBA layers rather
than a single raster output \cite{qwenlayered}. Its official examples emphasize
variable-layer decomposition and editable structure.

Portraits are a natural domain for this idea. A simple portrait editor might want
separate layers for background, hair, and face. CelebAMask-HQ provides pixel-level
face parsing masks for 30,000 high-resolution celebrity face images, making it a
reasonable source of supervision for such a task \cite{celebamaskhq}. The core
question for this project was therefore:

\begin{quote}
\emph{Can we fine-tune Qwen-Image-Layered with LoRA so that it decomposes portraits
into stable, editable, semantically meaningful RGBA layers?}
\end{quote}

The answer we found is mixed. The infrastructure can be made to work, and small
numeric improvements are possible, but straightforward LoRA fine-tuning did not
solve layer separation.

\section{Dataset and Target Construction}

\subsection{CelebAMask-HQ}

CelebAMask-HQ contains 30,000 face images selected from CelebA-HQ, with manually
annotated 512$\times$512 segmentation masks over facial components and accessories
such as skin, hair, eyes, eyebrows, mouth, lips, glasses, earrings, neck, and
clothing \cite{celebamaskhq}. This annotation structure makes it useful for face
parsing and facial image manipulation, but it also creates an important modeling
tension: segmentation masks are crisp semantic regions, while a generative layered
model may prefer broader visual chunks with soft alpha boundaries.

The raw dataset copy used in the project was not perfectly complete. Some image
IDs and mask IDs did not align. The preprocessing pipeline was therefore changed
to operate on the intersection of valid image IDs and valid mask IDs, recording
missing files instead of failing. This turned incomplete local/Drive data into a
usable training source.

\subsection{Layer Export}

The preprocessing pipeline exports each sample into:

\begin{itemize}[leftmargin=*]
    \item \texttt{layered\_composites/}: the input portrait composite,
    \item \texttt{layered\_layers/<sample\_id>/}: target RGBA layers,
    \item \texttt{metadata/layered\_samples.jsonl}: per-sample metadata.
\end{itemize}

For a source image $I(x,y)\in[0,1]^3$ and a binary mask $M_k(x,y)$ for target
layer $k$, the exported target layer is represented as an RGBA image
\begin{equation}
L_k(x,y)=
\left[
I(x,y)\odot M_k(x,y),\;
M_k(x,y)
\right],
\end{equation}
where the RGB channels are premultiplied by the alpha mask. This premultiplication
became important later because transparent regions should not carry arbitrary
hidden RGB values.

The first major target representation was a five-layer portrait scheme:
\begin{equation}
\mathcal{T}_5=
\{\mathrm{BG},\mathrm{HAIR},\mathrm{FACE\_SKIN},\mathrm{EYES},\mathrm{MOUTH}\}.
\end{equation}

Later, after observing weak visual improvement, we switched to a broader
three-layer scheme:
\begin{equation}
\mathcal{T}_3=
\{\mathrm{BG},\mathrm{HAIR},\mathrm{FACE}\},
\end{equation}
where
\begin{equation}
M_{\mathrm{FACE}}
=
M_{\mathrm{FACE\_SKIN}}
\lor
M_{\mathrm{EYES}}
\lor
M_{\mathrm{MOUTH}}.
\end{equation}
The background mask was defined as the complement of the foreground portrait
groups:
\begin{equation}
M_{\mathrm{BG}}
=
1-
\left(M_{\mathrm{HAIR}}\lor M_{\mathrm{FACE}}\right).
\end{equation}

\section{Colab and Packaging Infrastructure}

The project was run with a Drive-first Colab workflow. Raw data and processed
datasets lived on Google Drive, while smaller packaged artifacts were copied to
local \texttt{/content} storage for GPU training. This separation mattered because
Drive-backed random I/O over many small image files was too slow and fragile for
training.

The main Drive layout was:

\begin{verbatim}
/content/drive/MyDrive/CelebMaskHQ_Colab/
  repo/
  raw/
  processed/
  runs/
\end{verbatim}

Preprocessing and packaging were handled in CPU notebooks. Training, checkpoint
comparison, and inference evaluation were handled in GPU notebooks. We used a
paid Colab runtime exposing a G4-class GPU. The runtime reported a Blackwell-class
NVIDIA RTX PRO 6000 GPU with approximately 96 GB of VRAM; this matches Google
Cloud's documentation for G4 instances with NVIDIA RTX PRO 6000 Blackwell Server
Edition GPUs and 96 GB GPU memory per GPU \cite{googleg4}.

This runtime choice was also practical for Colab resource budgeting. In our
Colab account, the G4 runtime used a substantially smaller compute-unit rate than
the H100 option while still providing enough memory for Qwen-Image-Layered
inference, LoRA training, and checkpoint comparison. The resource snapshots in
Figure~\ref{fig:colab_usage} show the training and inference/evaluation contexts
used for the final experiments.

\begin{figure}[t]
\centering
\begin{minipage}{0.48\linewidth}
\centering

\includegraphics[width=\linewidth]{\detokenize{results/training_usage.png}}
\small Training runtime
\end{minipage}\hfill
\begin{minipage}{0.48\linewidth}
\centering

\includegraphics[width=\linewidth]{\detokenize{results/inference_usage.png}}
\small Inference/evaluation runtime
\end{minipage}
\caption{Colab resource snapshots for the G4-class runtime used in the
GPU stages. The setup offered enough memory for the project while keeping the
resource cost below the H100 runtime option available in Colab.}
\label{fig:colab_usage}
\end{figure}

Several pieces of infrastructure were needed before meaningful modeling
experiments were possible:

\begin{itemize}[leftmargin=*]
    \item metadata recovery from partially processed datasets,
    \item train/val sharded tar packaging,
    \item separate held-out test evaluation packaging,
    \item checkpoint save/resume logic with optimizer and scheduler state,
    \item LoRA adapter reactivation after resume,
    \item freshness checks to avoid stale Colab copies of Drive code.
\end{itemize}

This engineering work was not cosmetic. Without it, many failures looked like
modeling failures but were actually stale-code, packaging, or resume issues.

\section{Training Objective}

\subsection{Base Flow-Matching Loss}

The training script follows the Qwen layered pipeline's latent flow-matching
style objective. Let $z_0$ be the VAE latent of the target layered RGBA tensor and
let $\epsilon\sim\mathcal{N}(0,I)$ be Gaussian noise. For a sampled noise level
$\sigma$, the noisy target latent is
\begin{equation}
z_\sigma = (1-\sigma)z_0 + \sigma \epsilon.
\end{equation}

The transformer receives the packed noisy target latents together with packed
conditioning latents from the input portrait. Its prediction is trained against
the velocity-like target
\begin{equation}
v^\star = \epsilon - z_0.
\end{equation}

Ignoring masking, the original loss has the form
\begin{equation}
\mathcal{L}_{\mathrm{base}}
=
\mathbb{E}_{b,j}
\left[
w(\sigma_b)
\left(
\hat{v}_{\theta,bj}-v^\star_{bj}
\right)^2
\right],
\end{equation}
where $b$ indexes batch elements, $j$ indexes latent positions/channels/layers,
and $w(\sigma)$ is the scheduler/timestep weighting from the SD3-style training
utility.

\subsection{Why Plain MSE Was Problematic}

The target layers are sparse RGBA images. A face-feature layer or mouth layer
contains only a small visible region and a large transparent canvas. In early
experiments, plain latent MSE could be dominated by transparent regions or by
hidden RGB values in nominally transparent pixels. This was a plausible contributor
to two observed problems:

\begin{itemize}[leftmargin=*]
    \item stronger LoRA settings drifted into structured noise,
    \item weaker settings avoided collapse but barely changed the base model.
\end{itemize}

\subsection{Alpha-Aware Modification}

Two alpha-aware changes were introduced. First, target RGB values were premultiplied
by alpha before training export. Second, the loss was weighted by an approximate
latent-space visibility map derived from the target alpha channel.

Let $\alpha_k(x,y)$ be the target alpha mask for layer $k$. It is downsampled to
the latent grid with max pooling:
\begin{equation}
\bar{\alpha}_{kj}
=
\operatorname{MaxPool}(\alpha_k)_j.
\end{equation}

The latent loss weight is then
\begin{equation}
A_{kj}
=
\left(0.05 + 0.95\bar{\alpha}_{kj}\right)V_{kj},
\end{equation}
where $V_{kj}$ is a valid-layer mask. The small floor of $0.05$ preserves a weak
context signal for transparent regions while strongly emphasizing visible pixels.

The training loss used in the code is therefore
\begin{equation}
\mathcal{L}_{\alpha}
=
\frac{1}{B}
\sum_{b=1}^{B}
\operatorname{mean}_{j}
\left[
w(\sigma_b)
A_{bj}
\left(
\hat{v}_{\theta,bj}-v^\star_{bj}
\right)^2
\right].
\end{equation}

For logging, a normalized alpha-weighted MSE was also computed:
\begin{equation}
\operatorname{MSE}_{\alpha}
=
\frac{
\sum_{b,j}
A_{bj}
\left(
\hat{v}_{\theta,bj}-v^\star_{bj}
\right)^2
}{
\max\left(\sum_{b,j}A_{bj},1\right)
}.
\end{equation}

This modification improved stability, but it did not fully solve layer separation.

\section{Experiment Timeline}

\subsection{Large 5-Layer Train/Val Run}

The first major run used the five-layer target scheme with relatively strong LoRA
settings:

\begin{itemize}[leftmargin=*]
    \item learning rate $10^{-4}$,
    \item LoRA rank 16,
    \item LoRA alpha 16,
    \item gradient accumulation 8,
    \item 2000 training steps,
    \item checkpoints every 200 steps.
\end{itemize}

The loss curve looked plausible, and validation initially suggested a useful
checkpoint around step 1200. However, held-out qualitative comparison showed badly
degraded outputs with colorful structured noise. This was the first major lesson:
numeric training curves alone were not enough. The trained LoRA had damaged the
base model's useful decomposition behavior.

\subsection{Stability Rerun}

The next run reduced the learning rate and step count:

\begin{itemize}[leftmargin=*]
    \item learning rate $5\times10^{-5}$,
    \item LoRA rank 16,
    \item LoRA alpha 16,
    \item gradient accumulation 4,
    \item 800 steps,
    \item checkpoints every 100 steps.
\end{itemize}

The best validation checkpoint appeared much earlier, around step 200. Later
checkpoints generally worsened. The outputs were less catastrophically noisy than
the 2000-step run but still did not beat the base model convincingly.

\subsection{Weak Micro-Ablation}

To test whether the model needed only a very small update, we used:

\begin{itemize}[leftmargin=*]
    \item learning rate $10^{-5}$,
    \item LoRA rank 8,
    \item LoRA alpha 8,
    \item gradient accumulation 4,
    \item 100 steps,
    \item checkpoints every 25 steps.
\end{itemize}

This avoided collapse, but the outputs were visually close to the base model.
The adaptation was stable but nearly a no-op.

\subsection{Alpha-Weighted Micro-Ablation}

The alpha-aware changes were then added to the weak micro-ablation recipe. This
substantially improved the stability of the loss and helped avoid transparent-layer
pathologies. However, the visual outcome remained modest: the LoRA changed outputs
slightly but did not create reliably cleaner portrait layers.

\subsection{Target Representation Change: BG/HAIR/FACE}

The next hypothesis was that the five-layer target was too fine-grained and too
segmentation-like. The broader target representation was:

\begin{equation}
\mathrm{BG}/\mathrm{HAIR}/\mathrm{FACE}.
\end{equation}

This used a 5k processed subset to reduce preprocessing overhead and keep the
experiment focused on LoRA testing. The stable alpha-aware micro-ablation recipe
was reused:

\begin{itemize}[leftmargin=*]
    \item learning rate $10^{-5}$,
    \item rank 8,
    \item LoRA alpha 8,
    \item gradient accumulation 4,
    \item 100 steps,
    \item validation/checkpoint every 25 steps.
\end{itemize}

The best validation result was at step 50:

\begin{center}
\begin{tabular}{lcc}
\toprule
checkpoint & validation loss & note \\
\midrule
step 50 & 0.5089 & best checkpoint \\
step 75 & 0.8571 & worse \\
step 100 & 0.9682 & worse \\
step 25 & 1.1364 & worse \\
\bottomrule
\end{tabular}
\end{center}

Qualitatively, checkpoint 50 gave partial improvement on some samples, especially
by making one layer more hair-like or face-like, but the decomposition was still
not clean. The background layer often remained nearly full-canvas, and the face
layer frequently carried too much of the full portrait.

\subsection{Prompt Experiments}

We tested explicit layer-control prompts such as:

\begin{quote}
\texttt{decompose this portrait into exactly three editable RGBA layers: background, hair, and face}
\end{quote}

This hurt both the base model and the LoRA in qualitative comparisons. Later
inference-only prompt ablations found that the best behavior came from the
default call: no explicit prompt passed, with the pipeline's default English prompt
behavior enabled. Descriptive prompts such as ``a close-up portrait photo of a
person with visible face and hair'' often reduced layer quality or collapsed slots.

\subsection{Low-Learning-Rate Continuation}

We also tried continuing from checkpoint 50 with a lower learning rate and linear
decay. The goal was to refine the partially useful checkpoint without drifting.
Instead, validation became unstable:

\begin{center}
\begin{tabular}{lcc}
\toprule
checkpoint & validation loss & note \\
\midrule
step 55 & 0.5566 & best continuation checkpoint \\
step 60 & 1.2446 & worse \\
step 65 & 1.5236 & worse \\
step 70 & 3.2514 & severe drift \\
step 75 & 7.4991 & collapse-like behavior \\
\bottomrule
\end{tabular}
\end{center}

This reinforced the conclusion that simply training longer was not the right path.

\subsection{Three-Layer 200-Step Probe}

Because the earlier noisy failures came from the stronger five-layer setup, we
also tested whether the newer three-layer BG/HAIR/FACE representation would remain
stable at 200 steps. This run used the same 5k processed subset and the same
alpha-aware settings as the 100-step BG/HAIR/FACE run, but extended the maximum
step count to 200.

Unlike the old five-layer 2000-step run, the three-layer 200-step probe did not
collapse into colorful noise. However, it also did not produce a decisive
separation improvement. Its validation loss was still best around step 50:

\begin{center}
\begin{tabular}{lcc}
\toprule
validation step & validation loss & nearest saved checkpoint \\
\midrule
50 & 0.5087 & checkpoint 75 in the 200-step probe \\
75 & 0.8571 & checkpoint 75 \\
100 & 0.9680 & checkpoint 100 \\
200 & 0.7792 & checkpoint 200 \\
25 & 1.1369 & checkpoint 75 in the 200-step probe \\
\bottomrule
\end{tabular}
\end{center}

The checkpoint numbering detail matters: the 200-step probe saved checkpoints
starting at step 75 in that Colab run, so the original checkpoint-50 used in the
final comparison came from the earlier 100-step BG/HAIR/FACE run. The comparison
therefore evaluates base model vs. original checkpoint 50 vs. new checkpoint 200.

\section{Evaluation Metrics}

The final numeric evaluation compared the base model, the BG/HAIR/FACE
checkpoint-50 LoRA, and the BG/HAIR/FACE checkpoint-200 probe on 32 held-out test
samples using the default-call inference setting.

For a predicted binary alpha mask $\hat{M}_k$ and target mask $M_k$, we used:

\begin{equation}
\operatorname{IoU}_k
=
\frac{
|\hat{M}_k\cap M_k|
}{
|\hat{M}_k\cup M_k|
},
\end{equation}

\begin{equation}
\operatorname{Dice}_k
=
\frac{
2|\hat{M}_k\cap M_k|
}{
|\hat{M}_k|+|M_k|
}.
\end{equation}

To reduce threshold brittleness, we also used a soft alpha IoU:
\begin{equation}
\operatorname{SoftIoU}_k
=
\frac{
\sum_j \min(\hat{\alpha}_{kj},\alpha_{kj})
}{
\sum_j \max(\hat{\alpha}_{kj},\alpha_{kj})
}.
\end{equation}

Alpha MAE was:
\begin{equation}
\operatorname{MAE}_{\alpha,k}
=
\frac{1}{N}
\sum_j
|\hat{\alpha}_{kj}-\alpha_{kj}|.
\end{equation}

We also tracked layer-health diagnostics:

\begin{itemize}[leftmargin=*]
    \item empty-layer rate,
    \item full-layer rate,
    \item multi-layer alpha overlap,
    \item composite reconstruction MAE/PSNR.
\end{itemize}

The composite reconstruction metrics are only secondary. A model can reconstruct
the input well by putting most of the image into every layer, so reconstruction
quality does not imply editable separation.

\section{Final Results}

\begin{table}[h]
\centering
\caption{32-sample held-out evaluation. Higher IoU is better; lower alpha MAE is better.}
\begin{tabular}{lrrrrrr}
\toprule
model & macro IoU & soft IoU & alpha MAE & BG IoU & HAIR IoU & FACE IoU \\
\midrule
base model & 0.2461 & 0.2469 & 0.5461 & 0.3229 & 0.1607 & 0.2547 \\
checkpoint-50 LoRA & 0.2547 & 0.2554 & 0.5425 & 0.3229 & 0.1610 & 0.2800 \\
checkpoint-200 LoRA & 0.2585 & 0.2592 & 0.5512 & 0.3229 & 0.1578 & 0.2947 \\
\bottomrule
\end{tabular}
\end{table}

\begin{table}[h]
\centering
\caption{Layer-health diagnostics. Lower is generally better for these quantities.}
\begin{tabular}{lrrr}
\toprule
model & empty-layer rate & full-layer rate & multi-layer overlap \\
\midrule
base model & 0.1458 & 0.5208 & 0.8108 \\
checkpoint-50 LoRA & 0.1563 & 0.5313 & 0.8224 \\
checkpoint-200 LoRA & 0.1250 & 0.5313 & 0.8530 \\
\bottomrule
\end{tabular}
\end{table}

The LoRA improves macro IoU only modestly. Checkpoint 50 improves macro IoU by
0.0085 absolute, and checkpoint 200 improves it by 0.0123 absolute relative to the
base model. The gain is concentrated in the FACE layer: FACE IoU rises from
0.2547 for the base model to 0.2800 at checkpoint 50 and 0.2947 at checkpoint 200.
BG is essentially unchanged, and HAIR does not improve. Moreover, the layer-health
diagnostics still show poor separation. Checkpoint 200 has the highest multi-layer
overlap, 0.8530, meaning that the slight FACE IoU gain comes with worse layer
separation.

\begin{figure}[h]
\centering
\includegraphics[width=\linewidth]{\detokenize{results/summary_metrics_200step.png}}
\caption{Aggregate metrics for base model, checkpoint-50 LoRA, and checkpoint-200 LoRA.}
\end{figure}

\begin{figure}[h]
\centering
\includegraphics[width=0.75\linewidth]{\detokenize{results/per_layer_iou_200step.png}}
\caption{Per-layer IoU. The main gain is concentrated in the FACE slot; HAIR remains weak.}
\end{figure}

\begin{figure}[h]
\centering
\includegraphics[width=0.75\linewidth]{\detokenize{results/layer_health_diagnostics_200step.png}}
\caption{Layer-health diagnostics remain poor. The 200-step checkpoint does not reduce overlap; it increases the multi-layer overlap rate.}
\end{figure}

\subsection{Qualitative Showcase}

The qualitative comparison uses the final three-way evaluation: base model,
checkpoint 50 from the earlier BG/HAIR/FACE run, and checkpoint 200 from the
200-step probe. This lets us show the updated conclusion directly: the latest
three-layer setup does not collapse into obvious noise at 200 steps, but it still
does not learn robust, clean portrait layer separation.

\begin{figure}[h]
\centering
\includegraphics[width=\linewidth]{\detokenize{results/contact_sheet_04337_200step.png}}
\caption{A sample where the 200-step checkpoint improves the FACE slot relative to checkpoint 50. The output is not noisy, but separation is still imperfect: the FACE layer absorbs much of the portrait and the HAIR layer remains incomplete.}
\end{figure}

\begin{figure}[h]
\centering
\includegraphics[width=\linewidth]{\detokenize{results/contact_sheet_00361_200step.png}}
\caption{A sample where checkpoint 200 is worse than checkpoint 50. The longer-trained checkpoint does not become colorful noise, but it moves content into the wrong slots and reduces the macro IoU for this sample. This illustrates why 200 steps was not a reliable solution even though it remained visually stable.}
\end{figure}

Per-sample inspection showed that the mean gain was not broad. Some samples
benefited from checkpoint 200, while others degraded relative to checkpoint 50.
This makes the result a weak improvement in target-mask alignment rather than a
robust model-quality win.

\section{Failure Cases and Lessons}

\subsection{More Steps Can Be Worse}

The 2000-step five-layer run is the clearest noise-collapse failure case. Its loss
curve did not obviously announce disaster, but visual comparison showed noisy,
unusable layers. This should be understood as a failure of the earlier stronger
five-layer setup, not as the behavior of the final three-layer setup. The newer
BG/HAIR/FACE 200-step probe was visually more stable, but it still did not produce
reliable separation and increased overlap. The practical lesson is therefore more
specific: longer training did not monotonically improve decomposition quality,
even when the target representation was simplified.

\subsection{Validation Loss Did Not Fully Predict Editability}

Validation loss was useful for rejecting some checkpoints, but it was not a
complete proxy for visual layer quality. A checkpoint could have a better loss and
still produce layers that were not meaningfully more editable than the base model.
This motivated held-out visual comparison and eventually alpha-mask metrics.

\subsection{Prompt Control Was Fragile}

Explicit prompts that named the target layers seemed intuitive, but they worsened
the decomposition. The base model behaved best when allowed to use its default
prompting behavior. This suggests that Qwen-Image-Layered's prompt interface is
not a simple semantic slot controller for this portrait task.

\subsection{Semantic Segmentation Is Not the Same as Layer Decomposition}

CelebAMask-HQ supervision is semantic and crisp. Layer decomposition for editing
may be softer, broader, and more object-centric. A model trained for inherent
editability may not naturally organize a face image into the same masks used by a
face parsing benchmark.

\section{Future Directions}

The next stage should not simply train the same setup longer. More promising
directions are:

\paragraph{Explicit alpha and separation losses.}
The most direct fix is to supervise alpha structure explicitly with Dice/BCE mask
losses, overlap penalties, and full-canvas penalties for non-background layers.
For example, an added overlap penalty could take the form
\begin{equation}
\mathcal{L}_{\mathrm{overlap}}
=
\sum_{k<\ell}
\frac{1}{N}
\sum_j
\hat{\alpha}_{kj}\hat{\alpha}_{\ell j},
\end{equation}
penalizing multiple layers being opaque at the same pixel.

\paragraph{Hybrid segmentation-guided cleanup.}
A practical system could combine Qwen-Image-Layered with a face parsing or
segmentation model. Qwen would provide generative layer content, while a portrait
segmentation prior would clean the alpha masks. This would likely improve numeric
separation quickly, but it would be a hybrid system rather than pure LoRA
adaptation.

\paragraph{More model-aligned target representations.}
Instead of BG/HAIR/FACE, future targets could be broader:
\[
\mathrm{BG}/\mathrm{FULL\_SUBJECT}/\mathrm{HAIR\_DETAIL}
\]
or
\[
\mathrm{BG}/\mathrm{HEAD\_AND\_HAIR}/\mathrm{FACE\_FEATURES}.
\]
These may better match the pretrained model's native decomposition behavior.

\paragraph{Cleaner curated supervision.}
A smaller curated subset of frontal, minimally occluded portraits may reveal
whether noisy or ambiguous masks are limiting learning. This would separate target
noise from architectural or objective mismatch.

\section{Conclusion}

This project began as an attempt to fine-tune Qwen-Image-Layered into a portrait
layer decomposition model. The final result is not a strong model-quality win, but
it is a useful experimental finding. We built a working preprocessing, packaging,
training, resume, comparison, and evaluation pipeline. We showed that aggressive
LoRA training on the original five-layer target can damage the model, that weak
alpha-aware LoRA training avoids that collapse, and that a broader BG/HAIR/FACE
target remains stable but gives only small numeric gains.
The remaining bottleneck is not basic engineering instability; it is the alignment
between segmentation-derived supervision and the pretrained model's native layered
representation.

\begin{thebibliography}{9}

\bibitem{qwenlayered}
Qwen Team.
\textit{Qwen-Image-Layered model card}.
\url{https://huggingface.co/Qwen/Qwen-Image-Layered}.

\bibitem{celebamaskhq}
Cheng-Han Lee, Ziwei Liu, Lingyun Wu, and Ping Luo.
\textit{CelebAMask-HQ Dataset}.
\url{https://mmlab.ie.cuhk.edu.hk/projects/CelebA/CelebAMask_HQ.html}.

\bibitem{maskgan}
Cheng-Han Lee, Ziwei Liu, Lingyun Wu, and Ping Luo.
\textit{MaskGAN: Towards Diverse and Interactive Facial Image Manipulation}.
CVPR, 2020.

\bibitem{googleg4}
Google Cloud.
\textit{Accelerator-optimized machine family: G4 machine series}.
\url{https://cloud.google.com/compute/docs/accelerator-optimized-machines}.

\bibitem{nvidiag4}
NVIDIA.
\textit{NVIDIA and Google Cloud Accelerate Enterprise AI and Industrial Digitalization}.
\url{https://blogs.nvidia.com/blog/nvidia-google-cloud-enterprise-ai-industrial-digitalization/}.

\end{thebibliography}

\end{document}
